% Standalone mathematical notes for integration by the parent agent.
% Source question: incoming DHR, sections/09-research.tex, first subsection.
\section{The next two-particle resonance: all overlap terms and its cutoff constant}
\label{res:sec:resonance}\label{res:section}

The incoming \emph{Determinantal Harmonic Resonance} report
\cite[Section 9, first proposed target]{DHR} proves the
first resonance and explicitly proposes the two-particle next resonance
as its initial higher-resonance target.  We resolve that target here.
The input Gauss summation and the determinant identity are classical;
the contribution relative to that report is a normally convergent
subtraction at the next resonance, including the product overlap, an
explicit hard-cutoff constant, and its exact comparison with the
spectral Laurent constant.  No historical priority claim is intended.

Let $\Re a,\Re c>0$, and put
\[
 b_0=a+c+1,\qquad
 w_n(u)=\frac{\Gamma(n+a)\Gamma(n+c)}
 {\Gamma(n+b_0+u)\Gamma(n+1)},\qquad n\geq0.
\]
The two-particle series
\[
 Z(u)=\sum_{0\leq m<n}(n-m)^2w_m(u)w_n(u)
\]
initially converges for $\Re u>1$.  The next resonance is $u=0$.
Use the abbreviations
\begin{align*}
 t&=ac+a+c, & A&=\frac1{ac}, & h&=\frac{t}{ac}=1+\frac1a+\frac1c,\\
 L(u)&=\frac{\Gamma(a+1)\Gamma(c+1)\Gamma(1+u)}
 {\Gamma(a+1+u)\Gamma(c+1+u)}, &
 A(u)&=\frac{L(u)}{ac},\\
 B&=\psi(1)-\psi(a+1)-\psi(c+1), &
 T&=\psi'(1)-\psi'(a+1)-\psi'(c+1).
\end{align*}
Thus $L(0)=1$, $L'(0)=B$, and $L''(0)=B^2+T$.
For an independent anchor $\Re\alpha>0$, define
\[
 \kappa_\alpha(u)=-t+
       \left(\alpha-a-c-\tfrac12\right)u-\tfrac12u^2.
\]
Powers of $n+\alpha$ use the analytic logarithm on the right
half-plane, real on the positive real axis. For $x>0$, write
$x^{-u}=\exp(-u\log x)$ with the real logarithm.

\begin{theorem}[Normally convergent higher-resonance subtraction]
\label{res:thm:normal}
For $\Re u>-1$, the following three sums converge locally normally,
jointly in $a,c,u,\alpha$ on the stated right-half-plane domains:
\begin{align}
 A(u)&=\sum_{n\geq0}w_n(u),\label{res:eq:A}\\
 B_\alpha(u)&=\sum_{n\geq0}
    \left[nw_n(u)-(n+\alpha)^{-1-u}\right],\label{res:eq:B}\\
 D_\alpha(u)&=\sum_{n\geq0}\left[n^2w_n(u)
    -(n+\alpha)^{-u}
    -\kappa_\alpha(u)(n+\alpha)^{-1-u}\right].\label{res:eq:D}
\end{align}
Every fixed mixed parameter derivative may be taken termwise.  They
have the exact representations
\begin{align}
 B_\alpha(u)&=\frac{L(u)}u-\zeta(1+u,\alpha),\label{res:eq:Bexact}\\
 D_\alpha(u)&=\frac{L(u)(u+t)}{u(u-1)}
    -\zeta(u,\alpha)
    -\kappa_\alpha(u)\zeta(1+u,\alpha).\label{res:eq:Dexact}
\end{align}
Every apparent singularity of these expressions in $\Re u>-1$ is
removable.  The meromorphic continuation of the two-particle sum is
\begin{align}
 Z(u)={}&A(u)\left[\zeta(u,\alpha)
       +\kappa_\alpha(u)\zeta(1+u,\alpha)+D_\alpha(u)\right]
       -\left[\zeta(1+u,\alpha)+B_\alpha(u)\right]^2\notag\\
 ={}&-\frac{L(u)^2(1+u/a)(1+u/c)}{u^2(1-u)}.
 \label{res:eq:Z}
\end{align}
In particular, the square of the subtracted harmonic moment is an
essential overlap term.
\end{theorem}

\begin{proof}
The uniform Gamma-ratio expansion on compact parameter sets
\cite[Section 5.11(iii)]{DLMF} gives
\[
 w_n(u)=n^{-2-u}\left[1+\frac{\beta(u)}n+O(n^{-2})\right],\qquad
 \beta(u)=-t-(a+c+\tfrac12)u-\tfrac12u^2.
\]
When $n$ is replaced by $n+\alpha$ in the leading power of $n^2w_n$,
the next coefficient becomes $\beta(u)+\alpha u=\kappa_\alpha(u)$.
Consequently each summand in \eqref{res:eq:B} and \eqref{res:eq:D}
is $O(n^{-2-\Re u})$, locally uniformly; the same estimate without
subtraction applies to \eqref{res:eq:A}.  On a compact subset of
$\Re u>-1$ these estimates have a summable common majorant.
Cauchy's integral formula on slightly larger compact neighborhoods
then gives the assertion for all fixed mixed derivatives.  In
particular, differentiating at $u=0$ is justified before any
meromorphic expression is evaluated.

For the region in which the relevant ordinary moments converge,
Gauss summation \cite[Equation 15.4.20]{DLMF} yields
\begin{align*}
 \sum_{n\geq0}w_n(u)&=A(u),\\
 \sum_{n\geq0}nw_n(u)&=\frac{L(u)}u,\\
 \sum_{n\geq0}n(n-1)w_n(u)
    &=\frac{(a+1)(c+1)L(u)}{u(u-1)}.
\end{align*}
These formulas follow either by shifting the index in the defining
${}_2F_1$ series, or by differentiating it once and twice in its
argument before setting that argument to one.  Adding the last two
moments gives $L(u)(u+t)/(u(u-1))$.  Subtracting the two Hurwitz model
series proves \eqref{res:eq:Bexact} and \eqref{res:eq:Dexact} first on
nonempty open convergence domains and then by analytic continuation.
The proved holomorphy of the left sides proves removability at both
$u=0$ and $u=1$ wherever they occur in the displayed right sides.

Finite Cauchy--Binet gives
$\sum_{0\leq m<n\leq N}(n-m)^2w_mw_n=M_{0,N}M_{2,N}-M_{1,N}^2$.
Pass to the limit for $\Re u>1$ and insert the preceding moment
formulas.  The numerator simplifies using
$u(u+t)-ac(u-1)=(u+a)(u+c)$, proving \eqref{res:eq:Z}.

There is also a fully summable two-index remainder, rather than only
a formal determinant of continued quantities.  Write $b_n(u)$ and
$d_n(u)$ for the summands in \eqref{res:eq:B} and \eqref{res:eq:D}.
Then
\[
 A(u)D_\alpha(u)-B_\alpha(u)^2
 =\sum_{m,n\geq0}\left[w_m(u)d_n(u)-b_m(u)b_n(u)\right]
\]
is absolutely and locally normally convergent, by the product of
the three $\ell^1$ majorants just proved.  The remaining terms in
\eqref{res:eq:Z} are the explicitly displayed model sums and their
cross terms.  Thus the complete subtraction, including overlaps,
has a normally convergent realization independent of expanding
the Gamma product.
\end{proof}

\begin{theorem}[Hard cutoff and spectral finite part]
\label{res:thm:cutoff}
Let
\[
 S_N=\sum_{0\leq m<n\leq N}(n-m)^2w_m(0)w_n(0),\qquad x=N+1.
\]
As $N\to\infty$, locally uniformly for $\Re a,\Re c>0$,
\begin{equation}
 \boxed{S_N=A x-(\log x+B)^2-h(\log x+B)-2-A
             +O\left(\frac{\log^2(x+1)}x\right).}
 \label{res:eq:cutoff}
\end{equation}
The Laurent expansion of the spectral continuation is
\begin{equation}
 Z(u)=-\frac1{u^2}-\frac{2B+h}u
       -\left[2B^2+T+2hB+h+A\right]+O(u).
 \label{res:eq:Laurent}
\end{equation}
If $C_{\rm cut}$ denotes the constant after subtraction of
$Ax-\log^2x-(2B+h)\log x$, and $C_{\rm sp}$ the Laurent constant,
then
\begin{align}
 C_{\rm cut}&=-B^2-hB-2-A,\label{res:eq:Ccut}\\
 C_{\rm sp}&=2C_{\rm cut}+4+A-h-T.\label{res:eq:compare}
\end{align}
This uses the stated unnormalized spectral deformation $w_n(u)$;
neither constant is being assigned to an unspecified regulator.
\end{theorem}

\begin{proof}
Use the anchor $\alpha=1$ in Theorem~\ref{res:thm:normal}.
The expansion of $L(u)/u$ at zero gives
$B_1(0)=B-\gamma$.  Since $\zeta(0)=-1/2$ and
$\kappa_1'(0)=1/2-a-c$, the constant term in
\eqref{res:eq:Dexact} gives
\[
 D_1(0)=-t(B-\gamma)-ac-1.
\]
Consequently the three finite moments satisfy
\begin{align*}
 M_{0,N}&=A-\frac1x+O(x^{-2}),\\
 M_{1,N}&=\log x+B+O(x^{-1}),\\
 M_{2,N}&=x-t\log x-tB-ac-1+O(x^{-1}).
\end{align*}
The first formula follows by summing the $n^{-2}$ Gamma-ratio
tail; the other two follow from the proved summable residuals and
$H_x=\log x+\gamma+O(x^{-1})$.  Multiplication in the cutoff
determinant gives \eqref{res:eq:cutoff}.  In particular, the product
of the $-1/x$ correction in $M_{0,N}$ with the $x$ divergence in
$M_{2,N}$ contributes $-1$ to the constant.  Omitting it would give
an incorrect finite part.

Finally $L(u)^2=1+2Bu+(2B^2+T)u^2+O(u^3)$ and
\[
 \frac{(1+u/a)(1+u/c)}{1-u}
   =1+hu+(h+A)u^2+O(u^3).
\]
Insert these into \eqref{res:eq:Z}; this proves
\eqref{res:eq:Laurent} and the comparison formula.
\end{proof}

\paragraph{Dependence on the spectral normalization.}
Multiplying each particle's weight by a holomorphic function
$\lambda(u)$ with $\lambda(0)=1$ leaves the hard cutoff at zero
unchanged.  Write $\lambda_j=\lambda^{(j)}(0)$.  Since the total
spectral sum is multiplied by $\lambda(u)^2$, its residue and
constant become, respectively,
\[
 -(2B+h)-2\lambda_1,\qquad
 C_{\rm sp}-2\lambda_1(2B+h)-\lambda_1^2-\lambda_2.
\]
Thus the second derivative of the normalization enters at this
double pole.  These expressions follow by multiplying the two
Taylor--Laurent series; they explain exactly why the first-resonance
rule, which has only one amplitude derivative, cannot be copied
unchanged to the present resonance.

\paragraph{A convergent renormalized sum.}
Let $K_N=S_N-S_{N-1}$, $S_{-1}=0$, and $d=B-\gamma$.
The following is an ordinary absolutely convergent identity:
\begin{equation}
 \sum_{m\geq1}\left[
 K_{m-1}-A+\frac{2H_m+2d+h}{m}-\frac1{m^2}\right]
   =-d^2-hd-2-A.\label{res:eq:mass-sum}
\end{equation}
Indeed the largest-index formula
$K_N=w_N(0)[N^2M_{0,N-1}-2NM_{1,N-1}+M_{2,N-1}]$
and one further Gamma-ratio term give a remainder
$O(\log(m+1)/m^2)$ in the displayed summand.  Its partial sum is
exactly $S_{m-1}-Am+H_m^2+(2d+h)H_m$, and
\eqref{res:eq:cutoff} gives the limit.  The $m=1$ term is included;
the two-particle mass there is zero.

\subsection{Uniform spectral deformation and every cutoff jet}
The previous calculation extends to a uniform formula which retains
the two different escape powers.  In particular, derivatives at the
resonance can be computed without differentiating a nonuniform
asymptotic assertion.

\begin{theorem}[Uniform near-resonance expansion]
\label{res:thm:uniform}
Let $0<\delta<1/3$, $x=N+1$, and
$S_N(u)=\sum_{0\leq m<n\leq N}(n-m)^2w_m(u)w_n(u)$.
Uniformly for $|u|\leq\delta$ and for $a,c$ on compact subsets of
their right half-planes,
\begin{align}
 S_N(u)={}&Z(u)+\frac{A(u)x^{1-u}}{1-u}\notag\\
 &+x^{-u}\left[\frac{2L(u)}{u^2}
              -\frac{A(u)\beta(u)}u-\frac{A(u)}2\right]
       -\frac{x^{-2u}}{u^2(1-u^2)}
       +O\left(x^{-1+2\delta}\log^2(x+1)\right),
 \label{res:eq:uniform}
\end{align}
where $\beta(u)=-t-(a+c+1/2)u-u^2/2$.
The full displayed expression has a removable singularity at $u=0$.
The error is holomorphic in $u$ near this closed disk; every fixed
derivative at $u=0$ tends to zero as $N\to\infty$.
\end{theorem}
\begin{proof}
The uniform Gamma-ratio expansion from
Theorem~\ref{res:thm:normal} and the first two Euler--Maclaurin
terms for the corresponding Hurwitz tails
\cite[Equation 25.11.43]{DLMF} give
\begin{align*}
 M_{0,N}(u)&=A(u)-\frac{x^{-1-u}}{1+u}
                   +O(x^{-2+\delta}),\\
 M_{1,N}(u)&=\frac{L(u)-x^{-u}}u+O(x^{-1+\delta}),\\
 M_{2,N}(u)&=\frac{L(u)(u+t)}{u(u-1)}
       +\frac{x^{1-u}}{1-u}
       -x^{-u}\left[\frac12+\frac{\beta(u)}u\right]
       +O(x^{-1+\delta}).
\end{align*}
For a direct justification of the analytic tails, subtract
$n^{-1-u}$ from $nw_n(u)$ and subtract
$n^{-u}+\beta(u)n^{-1-u}$ from $n^2w_n(u)$ before summing from
$n=x$ to infinity.  The residuals are uniformly summable for
$\Re u>-1$.  Apply the ordinary Hurwitz tail expansion to the
model powers, first in their convergence regions and then by
analytic continuation.  This proves both the error bounds and
their uniform removability at $u=0$.

The apparent quotients involving $u^{-1}$ in the second and third
lines, after their pole cancellations, are
$O(x^\delta\log(x+1))$.  On multiplying the finite moment
determinant, the tail of $M_{0,N}$ times the leading divergence of
$M_{2,N}$ contributes
$-x^{-2u}/((1+u)(1-u))$.  All other terms involving a remainder
or this tail are $O(x^{-1+2\delta}\log^2(x+1))$.
Combine this contribution with the $-x^{-2u}/u^2$ from the square
of $M_{1,N}$; their sum is the last term in
\eqref{res:eq:uniform}.  The other terms simplify to the first
three displayed terms by \eqref{res:eq:Z}.
Finally, use Cauchy's estimate on a circle of radius $\delta$
around zero to transfer the tending-to-zero error to every fixed
derivative.  For clarity, the possible negative Laurent powers in
the full model cancel as follows: their $u^{-2}$ coefficients are
$-1+2-1=0$, and their $u^{-1}$ coefficients are
$-(2B+h)+(2B+h-2\log x)+2\log x=0$.
All other potential poles of the model lie outside $|u|<1$.
Thus the difference between the finite cutoff and the model is
holomorphic throughout $|u|<1$, as required for Cauchy's estimate.
\end{proof}

In particular, for each integer $q\geq0$ there are explicit
polynomials $P_q,Q_q$ of degrees at most $q,q+2$ such that
\begin{equation}
 \left.\partial_u^q S_N(u)\right|_{u=0}
       =xP_q(\log x)+Q_q(\log x)+o(1).
 \label{res:eq:all-cutoff-jets}
\end{equation}
They are obtained by taking $q!$ times the coefficient of $u^q$
in \eqref{res:eq:uniform}; only finitely many Gamma derivatives
are needed at each $q$.  The highest coefficient of the second
polynomial is universal:
\begin{equation}
 [y^{q+2}]Q_q(y)
   =\frac{(-1)^q(2-2^{q+2})}{(q+1)(q+2)}.
 \label{res:eq:universal-log}
\end{equation}
Indeed the only contributions of this degree come from
$2e^{-uy}/u^2-e^{-2uy}/u^2$.  This formula includes the
$-\log^2x$ at $q=0$ and the $+\log^3x$ at $q=1$.
The scale $x^{-2u}$ is essential even when only the finite
part of the undeformed cutoff is wanted.

\subsection{Generalized Stieltjes jets of the subtracted moments}
The normally convergent formula has an auxiliary Hurwitz anchor,
which supplies genuine generalized Stieltjes identities.  This is
a consequence of the subtraction theorem, not a claim that merely
differentiating Gauss summation is new.

Let $Y_q$ be the exponential complete Bell polynomial, with $Y_0=1$,
and set
\[
 \ell_j=\psi^{(j-1)}(1)-\psi^{(j-1)}(a+1)
                    -\psi^{(j-1)}(c+1),\qquad j\geq1.
\]
For $q\geq0$, write
\[
 W_q(n)=w_n(0)Y_q\bigl(-\psi(n+b_0),-\psi'(n+b_0),
                         \ldots,-\psi^{(q-1)}(n+b_0)\bigr).
\]
The interpretation is $W_0(n)=w_n(0)$.  Then
\begin{equation}
 \boxed{\sum_{n\geq0}\left[
 nW_q(n)-\frac{(-\log(n+\alpha))^q}{n+\alpha}\right]
 =\frac{Y_{q+1}(\ell_1,\ldots,\ell_{q+1})}{q+1}
       -(-1)^q\gamma_q(\alpha).}\label{res:eq:Stieltjes}
\end{equation}
All logarithms have the right-half-plane branch.  This is ordinary
absolute convergence, including every fixed mixed derivative in
$a,c,\alpha$.  For example,
\begin{align*}
 \sum_{n\geq0}\left[nw_n(0)-\frac1{n+\alpha}\right]
   &=B+\psi(\alpha),\\
 \sum_{n\geq0}\left[-nw_n(0)\psi(n+b_0)
                 +\frac{\log(n+\alpha)}{n+\alpha}\right]
   &=\tfrac12(B^2+T)+\gamma_1(\alpha).
\end{align*}
These follow from \eqref{res:eq:Bexact} and the defining expansion
$\zeta(1+u,\alpha)=u^{-1}+\sum_{q\geq0}(-1)^q\gamma_q(\alpha)u^q/q!$.
For the second subtracted moment one may use the similarly explicit
coefficient formula
\begin{align}
 \frac{D_\alpha^{(q)}(0)}{q!}
 ={}&-t p_{q+1}-(t+1)\sum_{j=0}^q p_j
     -\frac{\partial_s^q\zeta(0,\alpha)}{q!}\notag\\
 &-\kappa_{q+1}
     -\sum_{j=0}^{\min(2,q)}\kappa_j
                    \frac{(-1)^{q-j}\gamma_{q-j}(\alpha)}{(q-j)!},
 \label{res:eq:Djets}
\end{align}
where $p_j=Y_j(\ell_1,\ldots,\ell_j)/j!$,
$\kappa_0=-t$, $\kappa_1=\alpha-a-c-1/2$,
$\kappa_2=-1/2$, and $\kappa_j=0$ for $j>2$.
The left side is obtained by differentiating the summands of
\eqref{res:eq:D}; thus \eqref{res:eq:Djets} is a finite compiler
for convergent Gamma/polygamma sums in generalized Stieltjes
constants and ordinary spectral Hurwitz-zeta derivatives.

\subsection{A rational specialization and its complete polylogarithm generator}
At $a=c=1$, the weight is $w_n(0)=1/((n+1)(n+2))$.
With $x=N+1$ the elementary finite moments are
\[
 M_{0,N}=\frac{x}{x+1},\quad
 M_{1,N}=H_x+\frac2{x+1}-2,\quad
 M_{2,N}=x-3H_x-\frac4{x+1}+4.
\]
Their determinant simplifies to
\begin{equation}
 \boxed{S_N=x-H_x^2+H_x-1-\frac{H_x-1}{x+1},\qquad
 K_N=1-\frac{H_x}{x}-\frac{H_x-1}{x+1}.}\label{res:eq:rational}
\end{equation}
Thus the cutoff has harmonic degree two, but its largest-index
generator has depth two.  The finite difference removes the constant
coefficient of $H_x^2$; a formal harmonic-degree bound at the cutoff
level need not be sharp for the generator.

For $|z|<1$, summation of \eqref{res:eq:rational} gives
\begin{equation}
 \boxed{\sum_{0\leq m<n}
   \frac{(n-m)^2z^{n+1}}{(m+1)(m+2)(n+1)(n+2)}
 =\frac{z}{1-z}-1-\operatorname{Li}_2(z)
  -\frac{1+z}{2z}\log^2(1-z)-\frac{\log(1-z)}z.}
 \label{res:eq:G}
\end{equation}
The singularities at $z=0$ cancel, and its first term is $z^2/12$.
To check the summation, use
\[
 \sum_{m\geq1}\frac{H_mz^m}{m}
   =\operatorname{Li}_2(z)+\tfrac12\log^2(1-z),\qquad
 \sum_{m\geq1}\frac{H_mz^m}{m+1}
   =\frac{\log^2(1-z)}{2z}.
\]
Both follow by integrating the elementary generating function
$\sum H_mz^m=-\log(1-z)/(1-z)$.
In particular,
\begin{equation}
 \sum_{0\leq m<n}\frac{(n-m)^2\,2^{-n-1}}
 {(m+1)(m+2)(n+1)(n+2)}
 =2\log2-\log^22-\frac{\pi^2}{12}.\label{res:eq:half}
\end{equation}

For any integer $q\geq2$, let
\[
 J_q=\sum_{0\leq m<n}\frac{(n-m)^2}
 {(m+1)(m+2)(n+1)^{q+1}(n+2)}.
\]
It converges absolutely since $K_N\to1$.  We use the strict convention
$\zeta(j,1)=\sum_{m>k\geq1}m^{-j}k^{-1}$.  Put
$E_j=\zeta(j,1)+\zeta(j+1)$ for $j\geq2$,
$R_1=\zeta(2)-1$, and
$R_q=E_q-\zeta(q)-R_{q-1}$ for $q\geq2$.  Then
\begin{equation}
 J_q=\zeta(q)-E_{q+1}-R_q.\label{res:eq:Jall}
\end{equation}
Writing $G(z)$ for \eqref{res:eq:G} and $\Theta=z\,d/dz$, the
corresponding anchored primitives on $|z|<1$ are
\[
 \Theta^{-q}G(z)=\sum_{m\geq1}\frac{K_{m-1}z^m}{m^q}
  =\frac1{(q-1)!}\int_0^z\log^{q-1}(z/t)G(t)\frac{dt}{t}
       \quad(q\geq1),
\]
initially on a positive real segment and then by the analytic power
series.  Their endpoints at one are $J_q$ for $q\geq2$.
Indeed $R_q=\sum_{m\geq1}(H_m-1)/(m^q(m+1))$, and partial
fractions prove its displayed recurrence; the initial value follows
by telescoping.  Euler's classical reduction of $\zeta(j,1)$,
proved in the incoming report \cite[Section 6]{DHR}, gives
a terminating formula using only ordinary zeta values and products
of two such values, with weight at most $q+2$.
Explicitly, for $j\ge2$ that reduction is
\[
 \zeta(j,1)=\frac j2\zeta(j+1)
       -\frac12\sum_{r=1}^{j-2}\zeta(j-r)\zeta(r+1),
\]
with the empty sum interpreted as zero. Hence
\eqref{res:eq:Jall} is a finite ordinary-zeta evaluation at every
integer $q\ge2$.
The first two evaluations are
\begin{align}
 \boxed{J_2=3\zeta(2)-2\zeta(3)-\tfrac54\zeta(4)-1,}
 \label{res:eq:J2}\\
 J_3=1-2\zeta(2)+4\zeta(3)-\tfrac54\zeta(4)
                         -3\zeta(5)+\zeta(2)\zeta(3).
 \label{res:eq:J3}
\end{align}
The primitive $q=1$ diverges logarithmically and is deliberately
excluded from this endpoint statement.

\paragraph{Next questions.}
The next unsolved step after this result is an arbitrary-rank
subtraction at $b-a-c=2r-3$, where a two-by-two divergent moment
block interacts with a nontrivial lower determinant.  A second
target is the full family at $b-a-c=2r-2-M$: separate power and
logarithmic counterterms before comparing their cutoff constants
with spectral Laurent coefficients.  The example above also asks
for a depth invariant which respects finite differencing, rather
than only the degree of a formal cutoff polynomial.

% Primary reference URLs checked by this agent:
% https://dlmf.nist.gov/15.4.E20  (Gauss summation)
% https://dlmf.nist.gov/5.11.E13  (Gamma-ratio asymptotics)
% https://dlmf.nist.gov/25.11.E43 (Hurwitz asymptotics)
% The standalone report should cite these and the pinned incoming DHR report.
